%%%%%%%%%%%%%%%%%%%%%%%%%%%%%%%%%%%%%%%%%%%%%%%%%%
%%%%%%%%%%%%%%%%%%%%%%%%%%%%%%%%%%%%%%%%%%%%%%%%%%

% A primeira versão (1.0) deste modelo foi desenvolvida por Mauro Moura Severino, com a ajuda de Rubens Vasconcellos Terra Neto, para o Mestrado Profissional em Poder Legislativo (MPPL) da Câmara dos Deputados a partir do Modelo Canônico de TCC com ABNTEX2, elaborado pelo "abnTeX2 group". Cada uma das demais versões (2.0, 2.1 e 3.0) foi implementada por Mauro Moura Severino a partir da versão anterior.

%%%%%%%%%%%%%%%%%%%%%%%%%%%%%%%%%%%%%%%%%%%%%%%%%%

% Template de DISSERTAÇÃO
% Versão: v-3.0
% Data: 31/12/2024

% Esta versão viabiliza a produção de documentos com grau muito elevado de conformidade com as seguintes normas:
    % ABNT NBR 6023:2018 – Informação e documentação – Referências – Elaboração
    % ABNT NBR 6024:2012 – Informação e documentação – Numeração progressiva das seções de um documento – Apresentação
    % ABNT NBR 6027:2012 – Informação e documentação – Sumário – Apresentação
    % ABNT NBR 6028:2021 – Informação e documentação – Resumo, resenha e recensão – Apresentação
    % ABNT NBR 6034:2004 – Informação e documentação – Índice – Apresentação
    % ABNT NBR 10520:2023 – Informação e documentação – Citações em documentos – Apresentação
    % ABNT NBR 14724:2024 – Informação e documentação – Trabalhos acadêmicos – Apresentação
    % IBGE. Normas de apresentação tabular. 3. ed. Rio de Janeiro: IBGE, 1993.

%%%%%%%%%%%%%%%%%%%%%%%%%%%%%%%%%%%%%%%%%%%%%%%%%%
%%%%%%%%%%%%%%%%%%%%%%%%%%%%%%%%%%%%%%%%%%%%%%%%%%

% A EDIÇÃO DESTE ARQUIVO É DE RESPONSABILIDADE DO AUTOR, QUE DEVE EDITÁ-LO CONFORME INSTRUÇÕES A SEGUIR.

%%%%%%%%%%%%%%%%%%%%%%%%%%%%%%%%%%%%%%%%%%%%%%%%%%
%%%%%%%%%%%%%%%%%%%%%%%%%%%%%%%%%%%%%%%%%%%%%%%%%%

\acsetup{make-links=true} % Não edite este comando.

% ------------------------------------------------
% ------------------------------------------------
% Criação de siglas e abreviaturas
% ------------------------------------------------
% ------------------------------------------------

% ------------------------------------------------
% Para os comandos a seguir, substitua os conteúdos entre as chaves.
% ------------------------------------------------
% primeiras chaves: nome de chamada (apelido/label)
% terceiras chaves: sigla ou abreviatura
% quartas chaves: nome por extenso
% quintas chaves (se houver): nome por extenso no plural
% ------------------------------------------------

% SUGESTÃO DE USO: Não delete os comandos, copie e cole o comando que vai usar no espaço indicado a seguir, depois faça a edição.

%%%%%%%%%%%%%%%%%%%%%%%%%%%%%%%%%%%%%%%%%%%%%%%%%
% Espaço sugerido para a inclusão de siglas e abreviaturas.





%%%%%%%%%%%%%%%%%%%%%%%%%%%%%%%%%%%%%%%%%%%%%%%%%

% ------------------------------------------------
% Siglas no singular
% ------------------------------------------------

\DeclareAcronym{cd}{
  short = \normalfont{CD},
  long = {Câmara dos Deputados},
}

\DeclareAcronym{cefor}{
  short=\normalfont{Cefor},
  long={Centro de Formação, Treinamento e Aperfeiçoamento},
}

\DeclareAcronym{ppg}{
  short=\normalfont{PPG},
  long={Programa de Pós-Graduação},
}

\DeclareAcronym{mppl}{
  short=\normalfont{MPPL},
  long={Mestrado Profissional em Poder Legislativo},
}

\DeclareAcronym{usa}{
  short=\normalfont{USA},
  long={\emph{United States of America}},
}

\DeclareAcronym{ue}{
  short=\normalfont{UE},
  long={União Europeia},
}

\DeclareAcronym{urss}{
  short=\normalfont{URSS},
  long={União das Repúblicas Socialistas Soviéticas},
}

\DeclareAcronym{apf}{
  short=\normalfont{APF},
  long={administração pública federal},
}

\DeclareAcronym{ibama}{
  short=\normalfont{Ibama},
  long={Instituto Brasileiro do Meio Ambiente},
}

\DeclareAcronym{cpu}{
  short=\normalfont{CPU},
  long={\emph{central processing unit}},
}

% ------------------------------------------------
% Siglas no plural
% ------------------------------------------------

\DeclareAcronym{tcc}{
  short=\normalfont{TCC},
  long={trabalho de conclusão de curso},
  plural-form={trabalhos de conclusão de curso},
}

\DeclareAcronym{epi}{
  short=\normalfont{EPI},
  long={equipamento de proteção individual},
  plural-form={equipamentos de proteção individual},
}